\chapter{Implementation}\label{chap:implementation}

% NOTE (24 Nov 2026): Rewritten from the 25 October status log into a description of the system as built. Drafted with AI assistance at the owner's request; reword in your own words before submission. A drawn architecture diagram can replace Table 5.1 if wanted (the template has no TikZ; an image in figures/ works).

This chapter describes JadwAI as implemented at the end of the project. It covers the architecture, the data the administrator enters, how rules are entered and approved, the timetable solver, the independent validator, examinations, publication, and the JAI assistant with its locally hosted model. Each requirement's implementation status is listed in Appendix~\ref{app:requirements}, and screenshots of the interface are in Appendix~\ref{app:screenshots}. Measurements are reported in Chapter~\ref{evidence}.

\section{Architecture}\label{implementation-status}

JadwAI is a web application built for this project. It is not a fork of an existing timetabling system; UniTime~\cite{muller2025itc} served only as a reference. Table~\ref{tab:architecture} lists its parts.

\begin{table}[h]
\centering
\small
\begin{tabular}{@{}p{0.22\textwidth}p{0.27\textwidth}p{0.43\textwidth}@{}}
\toprule
Part & Technology & Role \\
\midrule
Web interface & React, TypeScript (Vite) & All administrator pages, the student and instructor portals, the JAI chat; English and Arabic with right-to-left layout \\
API server & Python, FastAPI & Authentication, permissions, data entry, rule review and approval, publication, exports, the JAI assistant \\
Database & PostgreSQL, SQLAlchemy, Alembic & Academic records, rules, saved inputs, timetables, audit events \\
Worker & Separate Python process & Runs generation jobs outside web requests \\
Solver & Google OR-Tools CP-SAT~\cite{google_cpsat} & Builds teaching and examination timetables \\
Validator & Separate Python module & Rechecks every timetable without using the solver's model \\
Language model & Ollama, Qwen2.5-7B-Instruct & Interprets typed rules and answers in the JAI chat, on the same server \\
\bottomrule
\end{tabular}
\caption{Parts of JadwAI.}
\label{tab:architecture}
\end{table}

The design follows one principle: every component that can make a mistake is checked by a component that does not share its logic. The language model or the rule parser only \emph{proposes} rules, and an administrator approves them. The solver only \emph{proposes} timetables, and the validator checks them. The browser carries out actions that JAI prepares only after the administrator confirms them, through the same endpoints the pages use, so every permission and validation check still applies.

\paragraph{Accounts and security.} There are administrator, viewer, instructor and student accounts.
\begin{itemize}
\item \textbf{Passwords and sessions.} Passwords are hashed with scrypt, and sessions expire. Sign-in attempts are rate-limited.
\item \textbf{Permissions.} Every request is checked on the server for the user's current permissions, and changing a user's permissions ends their sessions.
\item \textbf{Concurrent changes.} Changes carry a version number, so an outdated save is refused instead of overwriting a newer one.
\item \textbf{Repeated requests.} An operation repeated with the same request key is not applied twice.
\item \textbf{Audit.} Approvals, data changes, generation jobs and publications are recorded as audit events in the same database transaction as the change.
\end{itemize}

\section{Academic data}

The administrator enters the data in a fixed order: semester settings, rooms, instructors, student groups, then courses with their lectures, tutorials and laboratories. CSV import loads the same data with a preview and a list of errors, and it commits either everything or nothing.

\begin{itemize}
\item \textbf{Calendar.} Teaching days, the first and last hour, and named teaching periods. The demonstration calendar has six days and five 90-minute periods (08:30--10:00, 10:15--11:45, 12:00--13:30, 13:45--15:15, 15:45--17:15). A session may start at a period only if the consecutive periods from there cover exactly its length, so the breaks between periods are never used for teaching.
\item \textbf{Rooms} have a capacity, a type and features, for example a laboratory with GPUs. Rooms with the same capacity, type and features are treated as interchangeable, which reduces the solver's work.
\item \textbf{Instructors} are either doctors, who teach lectures, or teaching assistants (TAs), who teach tutorials and laboratories. Each course lists its eligible instructors separately for each kind of activity. Every TA must be given at least one session.
\item \textbf{Student groups} are organised as year, programme and group. Groups are the unit for student clashes: two sessions of one group never overlap.
\item \textbf{Advising students} are students whose courses differ from their group's, for example because of a retake. Each keeps the group's sessions except dropped courses and joins another group's session for each added course. Their personal sessions must not overlap, which the solver treats as a hard rule.
\end{itemize}

\section{Rules, interpretation and approval}

\paragraph{Rule types.} The supported rules are:
\begin{itemize}
\item instructor, room and group unavailability;
\item instructor preferred times;
\item permitted times for an activity or a group;
\item permitted or preferred times for all lectures, tutorials or laboratories;
\item a required room for an activity;
\item activities that must not share a day;
\item lectures before tutorials;
\item daily-hour and teaching-day limits for instructors and groups;
\item weekly teaching-hour limits for instructors;
\item locks that keep an instructor, time or room from an earlier timetable.
\end{itemize}
Each rule is mandatory or preferred. A mandatory rule holds in every valid timetable. A preferred rule adds a penalty to the convenience score when it is broken. Several permitted windows for one target are alternatives, while restrictions on different targets must all hold.

\paragraph{Daily policies.} Global policies limit students' and instructors' teaching hours per day and the number of consecutive periods, and can spread an activity's weekly sessions over different days. Teaching load is measured in elapsed teaching hours: a 90-minute session counts as 1.5 hours, once per instructor, however many groups attend.

\paragraph{Three ways to enter a rule.}
\begin{itemize}
\item \textbf{Forms.} Every rule type has a form.
\item \textbf{Built-in parser.} A typed statement is first read by a deterministic parser with patterns for English, formal Arabic and Egyptian Arabic. Arabic and English names are compared by their consonant sounds, so ``نور حسن'' matches the record ``Nour Hassan''. The parser never guesses: an unknown or ambiguous name, a missing day or an unrecognised request becomes a clarification code or an \emph{unsupported} code.
\item \textbf{Language model.} If the parser cannot interpret a statement and the language model is enabled, the statement goes to the model. The model receives the text, the calendar and only the records the text mentions. Ollama constrains its reply to the same JSON contract as the parser. The schema lists the proposals before the decision, so the model writes what it found before it decides whether the statement is ready (Section~\ref{baselines}). For statements about times for all classes, the model fills in a small semantic frame, and code does the calendar arithmetic.
\end{itemize}

\paragraph{Review and approval.} Every proposal, whatever its source, is validated against the current records and calendar. It is then shown beside the original statement with every field editable, and any clarification questions are listed. An approved proposal becomes a draft rule, and only rules the administrator then approves enter the inputs of the next generation. Approved rules can be revoked but not silently changed. Input changes never alter an existing timetable; instead, older drafts are marked as generated before the latest change.

\paragraph{External model adapter.} An adapter for a commercial model (GPT-5.4 mini) uses the same contract. It is switched off by default and was never called (Section~\ref{baselines}).

\section{Timetable generation}

\paragraph{Saved inputs.} Generation starts by saving an unchangeable copy of every input: calendar, records, approved rules and locks, objective profile and solver settings. The copy is stored with a fingerprint. The worker solves from this copy, so a timetable can always be traced to exactly the inputs it was built from.

\paragraph{Solver model.} Each session has separate decisions for its time, its room and, per activity, its instructor. This avoids a variable for every combination of time, room and instructor. An activity keeps one instructor for all its weekly sessions. No-overlap constraints keep the sessions of each room, instructor, group and advising student apart. Times are minutes from the start of the teaching week. The hard rules cover:
\begin{itemize}
\item complete assignment of every session;
\item fitting the teaching periods;
\item room capacity, type and features;
\item eligible instructors, with doctors for lectures and TAs for tutorials and laboratories;
\item every approved rule and lock;
\item weekly limits and the daily policies.
\end{itemize}

\paragraph{Objective.} The default objective, \texttt{convenience-v1}, minimises a weighted sum of penalties. In decreasing weight these are:
\begin{itemize}
\item sessions outside an instructor's preferred times;
\item repeated sessions of one activity on the same day;
\item empty periods between a group's first and last session of a day, and the same for instructors;
\item the number of days on which students and instructors must come to campus;
\item each group's busiest day;
\item sessions in the first or last period;
\item long teaching days;
\item uneven instructor workload.
\end{itemize}
Workload is penalised in rising steps, so that moving a session from a busy instructor to a quiet one is always worthwhile and teaching beyond the 12-hour standard load is a last resort. After solving, the worker recomputes every penalty from the timetable itself and checks the total against the solver's value.

\paragraph{Search settings and outcomes.} Default settings are 30~seconds and four search workers, or 240~seconds and eight workers for more than 150 weekly sessions. The solver first finds a valid timetable and then improves its score until the time limit. Its outcome is recorded as \emph{optimal}, \emph{feasible}, \emph{infeasible}, \emph{unknown} or \emph{model invalid}. A time limit reached without a timetable is reported as unknown, never as proof that no timetable exists. Jobs can be cancelled, and an interrupted job is recorded as failed with a reason.

\paragraph{Large timetables in parts.} One run places at most 600 sessions. A larger timetable is therefore built in parts of one or more programme years. Each part is solved around the fixed sessions of earlier parts, and every hard check runs on the combined timetable. Only a complete timetable can be published.

\paragraph{Input checks.} Before solving, conservative checks name problems that make a timetable impossible, for example a TA who cannot be given any session, an activity that no room can host, or more teaching hours than the suitable rooms provide. Passing these checks does not prove that a timetable exists.

\paragraph{Repair suggestions.} When no timetable exists, JadwAI looks for what the administrator could change: approved rules, locks, instructors' days off and daily limits. It relaxes all of them and then restores them a few at a time while a timetable still exists. What remains is a minimal set of changes, which is checked with the validator on a repaired timetable. The suggestions are only shown; they are never applied automatically. Physical limits, such as rooms, eligible instructors and the calendar, are never relaxed.

\section{Independent validation}

The validator is a separate module that does not import the solver's model-building code. It recomputes completeness, period fit, eligibility, the doctor/TA split, room suitability and capacity, overlaps for rooms, instructors, groups and advising students, approved rules, locks, weekly limits and daily policies. It works directly from the saved timetable and the saved inputs. Its report gives each violation with the sessions, resources, times and rules involved. It also lists how many items each check evaluated, which checks did not apply, and which limits were not configured. A check that could not run counts as not evaluated, never as passed. An incomplete or invalid result is never kept as a validated draft.

\section{Editing, comparison and publication}

Saved timetables are never changed. A manual move or a change of instructor creates a new draft from the same saved inputs, and the draft is kept only if the validator passes it. Two drafts can be compared session by session, and drafts can be named. Publication runs the validator again against the \emph{current} approved inputs, under an exclusive lock on the workspace. A database index guarantees that exactly one teaching timetable and one examination timetable are published at a time. The administrator can also publish one group's week from a draft: those groups' sessions, including shared lectures, are merged with the published timetable, and the merged result is validated before it replaces it.

Students and instructors see only the published timetable, and only their own sessions. Timetables are exported as server-generated landscape PDF pages with descriptive names, and as CSV. A statistics report gives sessions per day, overlap counts, room use and preference satisfaction.

\section{Examination timetables}

Each course has an exam plan: a number of quizzes, a midterm, an optional second-chance midterm and a final. Each exam type has its own length and start times. The number of students per exam is the sum of its groups' sizes, and each advising student is a cohort of their own. The examination model:
\begin{itemize}
\item assigns each exam to dates and halls whose seats cover its students, splitting a large exam over several halls when needed;
\item gives every hall in use a TA invigilator;
\item prevents any cohort, hall or invigilator from having overlapping exams.
\end{itemize}
It then minimises, in order of weight: extra exams for advising students, two exams on one day, exams on consecutive days, the busiest exam slot, and the number of halls in use. A lower bound on unavoidable consecutive-day exams is computed first. When the dates cannot avoid them, the administrator is told how many are unavoidable and what would reduce them, such as more exam dates or two exams per day.

\section{The JAI assistant}\label{evidence-jai}

JAI is a chat panel in JadwAI that answers questions and prepares tasks.

\paragraph{How a message is handled.}
\begin{enumerate}
\item \textbf{Built-in patterns.} Patterns in English, formal Arabic and Egyptian Arabic recognise common questions and commands immediately. Examples are whether a room or instructor is free, a group's timetable, when a course meets, free rooms, generating a timetable for a programme or year, exam timetables, PDFs, publishing and the state of the latest run. A request to publish the examination timetable is always answered by a pattern, because the fine-tuned model did not learn it (Section~\ref{evidence-jai-results}).
\item \textbf{Chat memory.} Each chat remembers the last room, instructor, group, course, days, time and timetable it discussed. A follow-up such as ``and on Monday?'' can be completed from this memory.
\item \textbf{The local model.} Messages that the patterns do not cover go to the local model. It may explain the facts supplied to it, discuss trade-offs and prepare at most one action. Every action is validated against the same contract as the corresponding page and is carried out only after the administrator confirms it.
\end{enumerate}
The model runs on the same server through Ollama, so no message or record leaves the machine. On the laptop used in this project, one model answer takes between one and six minutes. The measurements therefore concern behaviour, not interactive speed.

\paragraph{Verified guide.} A first evaluation showed that the model invented details about JadwAI, for example a ``third-party provider'' and instructors who ``can edit'' a published timetable. A guide of 25 short statements was therefore written, each checked against the application code. For each message, the (at most three) statements whose keywords match are included in the prompt. The instructions require the model to use only those statements for questions about JadwAI and otherwise to say that it is not sure.

\paragraph{Reply safeguards.} Before a reply is shown, the application:
\begin{itemize}
\item removes repeated sentences;
\item keeps the complete sentences of a reply that was cut off mid-loop;
\item rewrites a reply once into Egyptian Arabic when an Arabic message received an English reply;
\item cuts a reply at the first Chinese, Japanese or Korean character, a failure observed in the base model.
\end{itemize}

\paragraph{Mixed-direction text.} Arabic replies that contain English names, codes and times were displayed in a confusing order. For example, a time range of 08:30--10:00 inside an Arabic sentence appeared as 10:00--08:30. The chat now chooses each line's direction from its main language and isolates every English or numeric run, so such runs stay in reading order.

\paragraph{Fine-tuning.} Training examples are fictional conversations in English, formal Arabic and Egyptian Arabic, grouped into families (one intent with several wordings). Whole families go either to training or to the held-out set. The prompts are built by the same code as the live chat prompt, including the guide entries. From round~3 onwards, the dataset builder refuses an example whose answer relies on a guide entry that is missing from its prompt. Adapters were trained with QLoRA (rank~16) on Qwen2.5-7B-Instruct on a free T4 GPU, with the loss computed only on JAI's final answer.

\paragraph{Merging into Ollama.} Ollama does not load such adapters for this model, so each adapter was merged into Ollama's own 4-bit weights. A first merge back to the original 4-bit format erased about 98\% of the training, because the weight changes are smaller than the quantisation steps. Re-quantising each tensor to a different format (4-bit blocks to 6-bit, and 6-bit to 5-bit) preserved the training. A check over every trained matrix found on average 94--97\% of the trained change present per tensor type. Table~\ref{tab:jai-rounds} summarises the rounds, and their evaluation is reported in Section~\ref{evidence-jai-results}.

\begin{table}[h]
\centering
\small
\begin{tabular}{@{}lp{0.33\textwidth}rrrp{0.22\textwidth}@{}}
\toprule
Round & Training data & Rows & Epochs & Loss & Outcome \\
\midrule
v1 & Original examples, no guide in the prompts & 210 & 2 & 1.62 & Better behaviour, still invents facts \\
v2 & Same answers, guide added to the prompts & 210 & 3 & 1.11 & Worse; not used \\
v3 & Answers rewritten to follow the guide; ``not sure'' examples & 243 & 2 & 1.18 & Preferred in blind rating \\
v4 & v3 plus off-topic declines, topic changes, Arabic room terms, examination publishing & 317 & 2 & 1.07 & Preferred over v3 in blind rating; examination publishing still missed \\
\bottomrule
\end{tabular}
\caption{JAI fine-tuning rounds (training rows; held-out rows were 69, 69, 72 and 85).}
\label{tab:jai-rounds}
\end{table}

\section{Evaluation tooling}

The interpretation corpus, its annotation guide and the offline scorer are part of the repository (\texttt{evaluation/}). A manifest stores hashes of the data, the guide, the contract and the scoring code, so a changed file is detected. Default tools never read the reserved statements, and reading them needs a logged reason. The administrator interface has an evaluation workspace for human review, built but not yet used: it offers gold-free review downloads, in-browser annotation, immutable review records, comparison of two reviews with Cohen's $\kappa$, and an adjudication form. The fine-tuning tools (\texttt{finetuning/}) build the dataset from reviewed families, export it with hashes, merge adapters, check the merge, run the held-out evaluation and produce blind rating sheets.

\section{Testing}

The backend, the web interface and the evaluation tools each have an automated test suite, run before every reported measurement. Every fault found during the evaluation received a regression test. The current test counts are given in Chapter~\ref{evidence}.

\section{What is not implemented}

\begin{itemize}
\item Several terms at once, and a separate parent-course hierarchy.
\item Instructor aliases, and dated exceptions to availability.
\item Editable objective weights in the interface.
\item Scenarios created from a stated variation.
\item Minimal-change regeneration.
\item Calendar feeds and change notifications.
\item Instructors submitting their own availability.
\item The instrumentation for a participant study.
\item A complete audit export of model calls.
\item Formal accessibility and cross-browser checks.
\end{itemize}
Remote deployment and production operation were outside the project. The requirement-by-requirement status is in Appendix~\ref{app:requirements}.
